%%=============================================================================
%% Voorwoord
%%=============================================================================

\chapter*{\IfLanguageName{dutch}{Woord vooraf}{Preface}}%
\label{ch:voorwoord}

%% TODO:
%% Het voorwoord is het enige deel van de bachelorproef waar je vanuit je
%% eigen standpunt (``ik-vorm'') mag schrijven. Je kan hier bv. motiveren
%% waarom jij het onderwerp wil bespreken.
%% Vergeet ook niet te bedanken wie je geholpen/gesteund/... heeft


Binnen de opleiding Toegepaste Informatica aan HOGENT wordt in de laatste 2 jaren gewerkt met geautomatiseerde omgevingen om virtuele Windows- en Linuxsystemen klaar te zetten. Hierbij wordt gebruikgemaakt van Vagrant, een tool die het mogelijk maakt om labo- en ontwikkelomgevingen op een declaratieve manier aan te maken en provisionen. In de praktijk blijkt echter dat dit proces niet altijd even vlot verloopt. Opstarttijden kunnen traag zijn, provisioning kan vastlopen en bij onderbrekingen of foutmeldingen blijft het proces soms op de achtergrond doorlopen. Dit zorgt voor frustratie en tijdverlies bij studenten en docenten.
\\

Daarom heb ik voor deze bachelorproef gekozen om een eigen provisioner te ontwikkelen in Rust. De keuze voor dit onderwerp sluit aan bij mijn interesse in automatisering en programmeren, ook al volgt deze bachelorproef niet rechtstreeks uit mijn gekozen afstudeerrichting. Een provisioner combineert toch programmeren met systeembeheer en infrastructuur, waardoor het onderwerp goed aansluit bij de bredere richting van toegepaste informatica. Bovendien leek het interessant om na te gaan of een eigen oplossing performanter en gebruiksvriendelijker kan zijn dan de bestaande aanpak met Vagrant.


